\documentclass[10pt,a4paper]{article}
\usepackage[margin=1.65cm]{geometry}
\usepackage{mathptmx}
\usepackage[T1]{fontenc}
\usepackage[utf8]{inputenc}
\usepackage{booktabs}
\usepackage{array}
\usepackage{tabularx}
\usepackage{graphicx}
\usepackage[table]{xcolor}
\usepackage{microtype}
\usepackage{enumitem}
\usepackage{fancyhdr}
\usepackage{hyperref}
\usepackage{lastpage}
\usepackage{float}
\hypersetup{
  colorlinks=true,
  linkcolor=black,
  urlcolor=blue,
  pdftitle={QFlow P6 Final Same-Board FPGA Results},
  pdfauthor={Shahriar Rizvi},
  pdfsubject={Nexys3 H0-H4 physical comparison}
}
\definecolor{Navy}{RGB}{22,52,90}
\definecolor{Blue}{RGB}{40,95,150}
\definecolor{PaleBlue}{RGB}{235,243,250}
\definecolor{PaleGreen}{RGB}{235,248,239}
\definecolor{PaleGold}{RGB}{255,248,225}
\definecolor{PaleRed}{RGB}{253,238,238}
\definecolor{LightGray}{RGB}{245,245,245}
\setlength{\parindent}{0pt}
\setlength{\parskip}{4pt}
\setlist[itemize]{nosep,leftmargin=1.35em}
\renewcommand{\arraystretch}{1.16}
\pagestyle{fancy}
\fancyhf{}
\lhead{QFlow P6 Final Results}
\rhead{Nexys3 same-board comparison}
\cfoot{Page \thepage\ of \pageref{LastPage}}
\newcommand{\resultbox}[2]{
  \noindent\colorbox{#1}{\parbox{0.965\linewidth}{#2}}
}
\begin{document}

\begin{center}
{\Large\bfseries\color{Navy} QFlow P6 Final Same-Board FPGA Results}\\[4pt]
{\large H0--H4 Physical Comparison on Digilent Nexys3}\\[6pt]
Shahriar Rizvi \hfill Generated from frozen, hashed repository evidence
\end{center}

\resultbox{PaleGreen}{
\textbf{Final result. P6 completed: 10/10 steps.} Five independent H0--H4 bitstreams met the
100 MHz constraint and were cold-programmed into volatile FPGA SRAM.
Across 420 physical UART records, all 420 rows and all 3,780 result
fields exactly matched the frozen fixed-point model, with
\textbf{zero mismatches}.
}

\vspace{5pt}
\begin{tabularx}{\textwidth}{>{\bfseries}p{3.8cm}X}
Board / FPGA & Digilent Nexys3 Rev B / Spartan-6 XC6SLX16-2-CSG324 \\
Clock & 100 MHz; one kernel cycle = 10 ns \\
Physical modes & H0, H1, H2, H3 and H4 \\
Main proposed mode & H4 reinforcement-learned adaptive profile controller \\
H4 physical proof & 84/84 rows, 756/756 fields, display 4084, LD0 ON, LD1 OFF \\
Latency boundary & Kernel cycles only; UART, host capture, JTAG and display refresh excluded \\
\end{tabularx}

\section*{1. Same-board implementation and physical correctness}

\begin{table}[H]
\centering
\scriptsize
\resizebox{\textwidth}{!}{
\begin{tabular}{llrrrrrrr}
\toprule
Mode & Controller & Fmax & Cycles & Latency & Registers & LUTs & Slices & Exact \\
 & & (MHz) & (normal) & ($\mu$s) & & & & fields \\
\midrule
H0 & Shortest distance & 118.948 & 2.00 & 0.0200 & 493 & 801 & 301 & 756/756 \\
H1 & Fixed key-aware & 107.538 & 201.76 & 2.0176 & 1000 & 1484 & 589 & 756/756 \\
H2 & Fixed QFlow & 102.838 & 202.76 & 2.0276 & 1085 & 1650 & 649 & 756/756 \\
H3 & Rule adaptive & 102.648 & 202.76 & 2.0276 & 1097 & 1836 & 645 & 756/756 \\
\textbf{H4} & \textbf{RL adaptive} & \textbf{102.020} & \textbf{204.76} & \textbf{2.0476} & \textbf{1156} & \textbf{1706} & \textbf{677} & \textbf{756/756} \\
\bottomrule
\end{tabular}
}
\caption{Independent implementations using the same board, clock constraint,
replay order, fixed-point contract and physical result shell.}
\end{table}

\begin{minipage}[t]{0.49\textwidth}
\textbf{Achieved Fmax (all exceed 100 MHz)}\\[-2pt]
\begin{tabular}{ll}
H0 & \rule{5.055cm}{1.25ex} 118.948 MHz \\
H1 & \rule{4.570cm}{1.25ex} 107.538 MHz \\
H2 & \rule{4.371cm}{1.25ex} 102.838 MHz \\
H3 & \rule{4.363cm}{1.25ex} 102.648 MHz \\
H4 & \rule{4.336cm}{1.25ex} 102.020 MHz \\
\end{tabular}
\end{minipage}
\hfill
\begin{minipage}[t]{0.49\textwidth}
\textbf{Slice-LUT use}\\[-2pt]
\begin{tabular}{ll}
H0 & \rule{2.225cm}{1.25ex} 801 LUTs \\
H1 & \rule{4.122cm}{1.25ex} 1484 LUTs \\
H2 & \rule{4.583cm}{1.25ex} 1650 LUTs \\
H3 & \rule{5.100cm}{1.25ex} 1836 LUTs \\
H4 & \rule{4.739cm}{1.25ex} 1706 LUTs \\
\end{tabular}
\end{minipage}

\vspace{6pt}
\resultbox{PaleBlue}{
\textbf{H4 engineering trade-off versus H3.}
H4 uses \textbf{7.08\% fewer LUTs} than H3,
while registers change by 5.38\%, occupied
slices by 4.96\%, and normal kernel latency by
only 0.99\% (two cycles, or 20 ns).
Its Fmax changes by -0.61\% but remains above the
required 100 MHz.
}

\section*{2. Adaptive profile behavior}

\begin{table}[H]
\centering
\small
\begin{tabular}{llccccc}
\toprule
Mode & Controller & Profiles 0/1/2/3 & Transitions & Dwell runs &
Dwell range & Runtime Q update \\
\midrule
H3 & Threshold/rule adaptive & 24/36/18/3 & 44 & 45 & 1-7 & NO \\
H4 & Reinforcement-learned adaptive & 42/11/18/10 & 24 & 25 & 1-7 & NO \\
\bottomrule
\end{tabular}
\caption{The three intentionally invalid rows are excluded; each mode has
81 valid/no-path profile decisions.}
\end{table}

\resultbox{PaleGreen}{
\textbf{Main novelty demonstrated physically.}
H4 executed the frozen reinforcement-learned state-to-profile policy
using all four profiles (42/11/18/10), produced 24 profile transitions
and 25 dwell runs, and exactly matched the frozen model on every result
field. Relative to H3, H4 made 45.45\% fewer
profile transitions, demonstrating a more stable dwell-aware learned
profile sequence on this replay.
}

\section*{3. Held-out comparison and statistical interpretation}

\begin{table}[H]
\centering
\small
\resizebox{\textwidth}{!}{
\begin{tabular}{lrrrrrr}
\toprule
Comparison & Joint success & Mean diff. & Bootstrap 95\% CI &
Sign-test $p$ & Paired $d_z$ & Path changes \\
\midrule
H4 vs H2 & 58 & 86.47 & [-23.52, 228.22] & 0.219 & 0.176 & 6 \\
H4 vs H3 & 58 & -78.05 & [-242.53, 79.93] & 0.453 & -0.127 & 7 \\
\bottomrule
\end{tabular}
}
\caption{Paired bottleneck-fidelity comparison on the 64-row held-out
set. Positive differences favor H4.}
\end{table}

\resultbox{PaleGold}{
\textbf{Correct statistical claim.} The H4-H2 mean fidelity difference is positive, while the H4-H3 difference is negative. In both comparisons the bootstrap 95\% confidence interval crosses zero and the exact sign-test $p$-value exceeds 0.05. Therefore this replay supports physical correctness and adaptive execution, but not a statistically conclusive universal fidelity-superiority claim.
}

The held-out success rate is 90.625\% for every mode when the two
intentionally invalid rows are included. Raw scalar score is not
compared across modes because adaptive modes can use different profile
weights; comparing those scores directly would be invalid.

\section*{4. Thesis-ready contribution statement}

\begin{enumerate}[leftmargin=1.5em]
\item \textbf{Fair hardware comparison:} H0--H4 were implemented and
      tested independently on the same Nexys3 FPGA rather than comparing
      FPGA latency against software-only literature results.
\item \textbf{Bit-exact physical validation:} 420/420 physical rows and
      3,780/3,780 fields matched the frozen fixed-point golden model.
\item \textbf{Adaptive hardware behavior:} H3 proved threshold/rule
      switching; H4 proved reinforcement-learned adaptive profile
      selection with frozen dwell.
\item \textbf{Deterministic measurement:} all modes used the same
      100 MHz clock, vector order, record format and cycle boundary.
\item \textbf{Efficient realization:} H4 met timing with 1,156
      registers, 1,706 LUTs, 677 occupied slices, and no BRAM or DSP.
\end{enumerate}

\section*{5. Claim boundary and limitations}

\begin{table}[H]
\centering
\small
\begin{tabularx}{\textwidth}{>{\bfseries}p{4.2cm}X p{3.2cm}}
\toprule
Topic & Evidence-based statement & Status \\
\midrule
Physical correctness & 420 rows and 3,780 fields exactly match the frozen model & Supported \\
100 MHz operation & All five implementations have zero timing error and Fmax above 100 MHz & Supported \\
Learned adaptation & H4 follows a frozen learned state-to-profile and dwell policy & Supported \\
Online learning & No runtime Q-table update occurs on the FPGA & Do not claim \\
Universal fidelity superiority & Held-out confidence intervals include zero & Not established \\
Power / energy superiority & No common activity-calibrated H0--H4 evidence & Not measured \\
Photo/video evidence & OMITTED\_NOT\_ARCHIVED\_IN\_REPOSITORY & Explicitly labelled \\
\bottomrule
\end{tabularx}
\end{table}

\resultbox{PaleRed}{
\textbf{Defense-safe wording.}
The evidence supports a reinforcement-learned adaptive FPGA controller
that is physically correct, deterministic, timing-clean and resource
feasible. It does not support online learning, universal network
superiority, or power superiority.
}

\section*{6. Traceability}

\begin{tabularx}{\textwidth}{>{\ttfamily\scriptsize}p{5.2cm}X}
same\_board\_comparison.csv & Timing, cycles, resources and correctness \\
heldout\_pairwise\_h4\_statistics.csv & Confidence intervals, sign tests and effect sizes \\
profile\_switching\_dwell.csv & Profile transitions and dwell \\
physical\_evidence\_matrix.csv & Programming, UART hashes, displays and LEDs \\
power\_energy\_status.csv & Explicit power/energy omission \\
\end{tabularx}

\end{document}
